\begin{workedexample}[Worked Example: Far-field distance]
A reflector of largest dimension $D = 0.3\ \text{m}$ is measured at
$f = 10\ \text{GHz}$ ($\lambda = 0.03\ \text{m}$). The far field begins at
\[ R \ge \frac{2D^2}{\lambda} = \frac{2(0.3)^2}{0.03} = 6\ \text{m}. \]
Measuring closer than this corrupts the pattern; if $6\ \text{m}$ of range is
impractical, a near-field scan with a far-field transform is used instead.
\end{workedexample}

\begin{keytakeaways}
\begin{itemize}[leftmargin=1.4em,itemsep=2pt]
\item Far field begins at $R \ge 2D^2/\lambda$; measure in an anechoic chamber or on a range.
\item Gain by comparison (to a standard) or the three-antenna method (no standard needed).
\item Near-field scanning + transform yields the far-field pattern when a far-field range is impractical.
\end{itemize}
\end{keytakeaways}

\begin{exercises}
\item Far-field distance for $D = 1\ \text{m}$ at $3\ \text{GHz}$?
\item Which gain method needs no pre-calibrated reference antenna?
\item How is antenna impedance/return loss measured?
\end{exercises}

\furtherreading{Balanis~\cite{balanis} (ch.~17); Pozar~\cite{pozar} (ch.~4).}
